% SPDX-License-Identifier: LPPL-1.3c
% Copyright (C) 2026 Christophe Jorssen
% Author and Current Maintainer: Christophe Jorssen <christophe.jorssen@gmail.com>
% LPPL maintenance status: maintained
% This file is part of luafemm. See LICENSE and LICENSES.md for its terms.
\section{Analytical gaps on reduced sectors}
\label{sec:sector-gaps}

A rotationally repeating machine can be represented by one sector. Pair its
radial boundaries, then connect its rotor and stator through an analytical air
gap. The two bulk meshes remain in their own reference frames for calculation.
For display, the rotor is rotated, continued periodically and clipped against
the stator sector. Both geometry and field pictures therefore use the physical
stator frame, including when a rotor feature crosses a radial seam. Scalar
potential changes sign on an antiperiodic copy; the magnetic vectors also rotate.
Profiles sample the same physical field. A drawing transformation still only
changes presentation, not the mechanical angle.

The sector must fit the declared domain and its rigid periodic pairing. All
physical regions stay clear of both gap traces. The gap contains source-free
vacuum air. Start with a coarse mesh and a modest segment count, then refine the
bulk mesh and both traces to check convergence of linkage and torque.

\luafemmexamplebegin{sector-gap}
\begin{codeexample}[width=0pt]
\begin{tikzpicture}[
  femm/problem={xmin=0,xmax=10,ymin=0,ymax=10,mesh size=.8,depth=10},
  femm/pair={first=radial-a,second=radial-b,sectors=4,type=antiperiodic},
  femm/sliding band={name=quarter-gap,method=analytical,
    inner radius=4,outer radius=5,segments=12,sectors=4,
    periodicity=antiperiodic},scale=2]
  \draw[femm/domain] (0,0)--(10,0) arc[start angle=0,end angle=90,radius=10]--cycle;
  \path[femm/frontier={name=radial-a,type=periodic}] (0,0)--(10,0);
  \path[femm/frontier={name=radial-b,type=periodic}] (0,0)--(0,10);
  \filldraw[fill=orange!25,
    femm/region={material=copper,current density=1e8}]
    (1,1) rectangle (2,2);
  \filldraw[fill=orange!25,
    femm/region={material=copper,current density=-1e8}]
    (7,1) rectangle (8,2);
  \draw[blue!50!green,dashed] (4,0) arc[start angle=0,end angle=90,radius=4];
  \draw[blue!50!green,dashed] (5,0) arc[start angle=0,end angle=90,radius=5];
\end{tikzpicture}
\end{codeexample}

The domain and the template are now complete. Choose a position whose inner
winding crosses a radial seam. The neighbouring antiperiodic copy enters the
visible sector with the opposite source sign. The field and the torque below
come from that same position.
\begin{codeexample}[]
\begin{tikzpicture}[
  femm/problem={sliding band=quarter-gap,rotor angle=53},scale=3]
  \pic {femm geometry};
  \pic[femm/lines=13] {femm field};
  \node[above] at (5,10) {$\theta=53^\circ$};
\end{tikzpicture}

\begin{tikzpicture}[
  femm/problem={sliding band=quarter-gap,rotor angle=53}]
  \femmmeasure[name=quarter-torque,quantity=force,method=band]
  \node {$T_{\rm sector}(53^\circ)=
    \femmvalue[component=torque]{quarter-torque}$};
\end{tikzpicture}
\end{codeexample}
\luafemmexampleend{sector-gap}

\subsection{Sector controls}
\begin{key}{/femm/sliding band/sectors=\meta{integer} (initially 1)}
One retains the full-circle interface. Values from 2 to 32 select a reduced
sector of angle $360^\circ/N$ and require |method=analytical|, an explicit
|femm/domain| and matching paired radial frontiers. |segments| then counts
segments per sector, from 4 to 128. Fourier mode resolution is still checked
against the equivalent full-circle traces; the retained maximum order is 256.
\end{key}
\begin{key}{/femm/sliding band/periodicity=\meta{choice} (initially periodic)}
|periodic| identifies the two sector traces with the same potential.
|antiperiodic| reverses its sign and requires an even sector count.
This choice must agree with the radial frontier pair. For $N$ sectors, periodic
traces contain harmonics $kN$; antiperiodic traces contain $(k+\tfrac12)N$.
Temporal source phase remains independent of this spatial sign convention.
\end{key}
\begin{key}{/femm/sliding band/start angle=\meta{degrees} (initially 0)}
Start direction of the sector in the reference coordinates, measured
counterclockwise about |center x|, |center y|. The domain and frontier pairing
must use the same centre and angular extent. Full-circle templates require zero.
\end{key}

Energy and transmitted torque are reported for the modelled sector. For torque
about the rotation axis, set the measurement keys |origin x| and |origin y| to
the gap centre; their default is the coordinate origin. Multiply
scalar energy and torque by |sectors| for a full machine with the specified
symmetry. Cartesian force components require rotating and summing each sector's
contribution; simply multiplying one sector's force vector is generally wrong.
The gap traction integral is not a closed contour around an isolated sector body.
Its torque is the torque transmitted across the circular interface.

The radial Laplace problem is solved per retained Fourier mode. Angular P1 trace
functions are integrated analytically, with signed periodic continuation and a
rotation of the inner trace. The weak operator is scaled by the sector fraction.
The same analytic curl is used for gap sampling and traction. Virtual work and
modal annular solutions provide independent checks of torque and normalisation.
Bulk meshes and signed pair constraints are reused at each angle and validated
on restoration. Eccentric motion is available through full-circle P1 templates, not through
this concentric operator. Network-driven moving templates remain unsupported.
Native AGE interpolation is a separate choice (Section~\ref{sec:native-gaps}).

Bulk contour pieces are transformed exactly and clipped to the displayed wedge.
Inside the gap, the analytic potential is sampled on a polar rendering grid;
|contour subdivisions| refines that grid without changing the solve.
The bulk P1 traces and the truncated Fourier reconstruction are different
interpolants. Their contour pieces retain separate endpoints, so the renderer
does not invent connections between unequal values. Refine the mesh and retained
modes to assess trace mismatch. Energy and torque always use the calculation
operator, never this rendering grid.

\subsection{Preserving native AGE answer files}
FEMM's AGE quad-point records and the Fourier interface above describe different
interpolations. Relabelling one as the other would silently change the numerical
model. Native answer records can instead be parsed and preserved without
recalculation, including FEMM 5's optional raw ring lists. For field import, explicitly select the generating solver's interpolation
with |age interpolation=femm4| or |femm5|; see Section~\ref{sec:native-gaps}.
Passive copying itself needs no such choice and performs no basis conversion.

\begin{command}{\femmcopyanswer\oarg{options}}
Read, validate and write a passive magnetic answer. It preserves its embedded
problem, nodal potentials, element records, circuit states, periodic pairs and
AGE quad records. No magnetic picture is required and no solve occurs. Input
and output filenames must differ. Relative outputs follow the engine output
directory; input lookup also searches that directory after the format resolver.
External previous-solution companions are not copied or rebased. Keep their
relative paths valid when copying an answer. The same operation is available through
|/tikz/femm/copy answer={options}|.
\end{command}
\begin{key}{/tikz/femm/copy answer=\marg{options}}
Equivalent to |\femmcopyanswer|; this is an explicit file-writing operation.
\end{key}
\begin{key}{/femm/copy answer/input=\meta{filename} (initially empty)}
Required source |.ans| filename. Supported fixed planar time records are
preserved without assigning an absolute clock or performing a solve. Unknown
trailing records, malformed counts, indices and nonfinite numbers are rejected.
Copying retains the companion-file guard at both paths.
\end{key}
\begin{key}{/femm/copy answer/output=\meta{filename} (initially empty)}
Required destination filename. Numeric data are written with 17 significant
digits and checked by the same bounded parser before writing.
\end{key}
\begin{key}{/femm/copy answer/rings=\meta{choice} (initially preserve)}
|preserve| retains FEMM 5's |[AgeRings]| lists when present. |omit| writes only the
common AGE records read by FEMM 4 and FEMM 5. Omitting raw rings loses the data
needed for FEMM 5 to re-pair that saved mesh at another angle. This option does
not translate material laws or certify that every other feature of a FEMM 5
problem is supported by FEMM 4.
\end{key}

The following executed example first saves an ordinary answer, then exercises
the same preservation command. For a native AGE file, use its filename as
|input|; no conversion of the AGE operator is performed.
\luafemmexamplebegin{answer-preservation}
\begin{codeexample}[]
\begin{tikzpicture}[
  femm/problem={xmin=-5,xmax=5,ymin=-5,ymax=5,
    mesh size=2,depth=10,interpolation=femm}]
  \filldraw[fill=red!20,femm/region={material=cast-alnico-5-lng37}]
    (-2,-1) rectangle (2,1);
  \pic[femm/lines=7] {femm field};
  \femmexport[file=\jobname-source.ans,format=ans]
\end{tikzpicture}
\femmcopyanswer[input=\jobname-source.ans,output=\jobname-copy.ans]
\end{codeexample}
\luafemmexampleend{answer-preservation}
